% i_entries.tex - Terms beginning with I
\chapter*{I}
\addcontentsline{toc}{chapter}{I}

\dictentry{ISO/IEC 14443}{/ˌaɪ ɛs ˈoʊ aɪ iː ˈsiː fɔːr-tiːn fɔːr-fɔːr θriː/, \textit{n.}}{
The international standard for contactless smart cards and proximity readers. It
defines the physical characteristics and communication protocols used by
contactless cards to exchange data with a terminal over short distances
(typically under 10 centimeters) at a frequency of 13.56 MHz.

\example{When you tap your credit card or smartphone at a checkout counter, the
transaction is powered by the ISO/IEC 14443 standard, which ensures the reader
can identify your card wirelessly.}

\analogy{Biological}{Like \textbf{Juxtacrine Signaling} in cells. Just as
certain biological signals require direct contact or extreme proximity between
cells to transmit information, ISO/IEC 14443 ensures that data is only
exchanged when the card is "touching" the reader's field. This proximity acts as
a physical security layer, preventing signals from being intercepted from a
distance.}

\crossref{\textbf{Contactless}, \textbf{EMV}, \textbf{NFC}, \textbf{RFID}}

\etymology{Named after the International Organization for Standardization (ISO)
and the International Electrotechnical Commission (IEC). The number 14443 is
part of their sequential numbering system for standards.}
}
\indexentry{ISO/IEC 14443}

\dictentry{ISO/IEC 15693}{/ˌaɪ ɛs ˈoʊ aɪ iː ˈsiː fɪfˈtiːn sɪks-naɪn-θriː/, \textit{n.}}{
The international standard for "vicinity" contactless smart cards and RFID tags.
Unlike proximity cards (ISO/IEC 14443), vicinity cards are designed to be read
from a much greater distance, typically up to 1.5 meters (5 feet). It operates
at the same 13.56 MHz frequency but is optimized for range over high-speed data
transmission.

\example{Ski passes often use ISO/IEC 15693 technology, allowing you to pass
through a gate while the card is still in your jacket pocket, without needing
to "tap" it directly against a reader.}

\analogy{Biological}{Like \textbf{Paracrine Signaling} in biology. While some
cells require physical contact to talk (Juxtacrine), paracrine signaling
involves a cell secreting a signal that travels through the local environment
to reach nearby neighbor cells without direct touching. Similarly, ISO/IEC 15693
allows a reader and card to "talk" across a small room or hallway, rather than
requiring a physical "tap."}

\crossref{\textbf{ISO/IEC 14443}, \textbf{NFC}, \textbf{RFID}, \textbf{Vicinity Card}}

\etymology{A sequential standard number from the ISO/IEC series, following the
14443 standard as contactless technology evolved to support greater distances.}
}
\indexentry{ISO/IEC 15693}

\dictentry{ISO 18000}{/ˌaɪ ɛs ˈoʊ eɪˈtiːn ˈθaʊzənd/, \textit{n.}}{
A family of international standards for Radio Frequency Identification (RFID)
used specifically for item management, logistics, and supply chain tracking. It
defines communication protocols across various frequencies—from Low Frequency
(LF) and High Frequency (HF) to Ultra-High Frequency (UHF) and Microwave—to
ensure items can be tracked globally as they move between different systems.

\example{A shipping company uses ISO 18000-compliant RFID tags on every pallet
in a warehouse. This allows automated scanners to log the movement of thousands
of items simultaneously without manual barcode scanning.}

\analogy{Biological}{Like \textbf{Pheromone Trails} in an ant colony. Ants use
specific chemical signals to label "items" (food) and create paths that the
entire colony can follow across vast distances. Just as different pheromones
serve different purposes (marking a path vs. signaling danger), ISO 18000
provides different "frequencies" of signals to manage items across the complex
environment of a global supply chain.}

\crossref{\textbf{ISO/IEC 14443}, \textbf{ISO/IEC 15693}, \textbf{RFID}, \textbf{Logistics}}

\etymology{A broad standard number designated for the "Automatic identification
and data capture techniques" category by ISO/IEC.}
}
\indexentry{ISO 18000}

% Additional I terms will be added here
